\documentclass[11pt]{article}
% Shared local declarations. Document roots use the repository's business-admin family.
\usepackage{business-admin}
\usepackage{amsmath,amssymb,tabularx,longtable,array,booktabs,enumitem}
\usepackage{ragged2e}
\newcolumntype{Y}{>{\RaggedRight\arraybackslash}X}
\newcommand{\SuiteMeta}[2]{\def\DocTitle{#1}\def\DocSubtitle{#2}\def\DocVersion{1.0}\def\DocStatus{Educational and operational reference}\def\DocOwner{Parliamentary Procedure and Governance Suite}\def\DocSummary{Independently authored reference and working materials, based on RONR 12th edition.}}
\newcommand{\SuiteOpen}{\makedoctitle\vspace{-0.5em}\noindent\rule{\linewidth}{0.4pt}\par\medskip}
\newcommand{\RonrAnchor}[1]{\par\smallskip{\footnotesize\textbf{Primary reading:} Robert's Rules of Order Newly Revised, 12th ed., #1. Original synthesis; not official interpretation.}\par}
\newcommand{\WriteLine}{\par\noindent\rule{\linewidth}{0.35pt}\par\vspace{0.5cm}}
\newcommand{\WriteLines}[1]{\par\noindent\rule{\linewidth}{0.35pt}\par\vspace{0.35cm}\noindent\rule{\linewidth}{0.35pt}\par\vspace{0.35cm}\noindent\rule{\linewidth}{0.35pt}\par\vspace{0.25cm}}
\newcommand{\Field}[1]{\par\medskip\noindent\textbf{#1}\par\WriteLines{3}}
\newcommand{\GovCheckbox}{\(\square\)}
\setlength{\emergencystretch}{3em}

\SuiteMeta{Motion Processing and Chair Rulings}{Technical reference | original application synthesis}
\begin{document}
\SuiteOpen
\section{Scope and interpretation}
This is an original operational reference based on RONR (12th ed.). It is subordinate to applicable law and the organization’s validly adopted governing documents. Rule citations identify the primary source; the authors of this suite have not issued official interpretations.

\section{Normal transaction}
Recognize an eligible member, receive a properly framed motion, obtain a second where needed, state the question, permit applicable debate and subsidiary motions, put the question, and announce the result. Ownership of a stated motion and withdrawal rules matter.

\textbf{Practical audit hook.} Source anchor: RONR 3--7, 10.

\section{Characteristics matrix}
For each motion inspect precedence, second, debatability, amendability, vote requirement and reconsiderability. These descriptors vary by context; do not run a rigid workflow that ignores exceptions.

\textbf{Practical audit hook.} Evidence: one row for every distinct pending motion, with its text and governing section.

\section{Escalating procedural disputes}
When a point of order is raised, identify the alleged procedural breach rather than policy merits. Record the ruling and basis; if an appeal is proper, frame the question around sustaining the chair’s decision. Do not confuse an appeal with a no-confidence vote.

\textbf{Practical audit hook.} Source anchor: RONR 23--24 and 62:2--7.

\section{Consent and fairness}
Unanimous consent can streamline routine action only when there is no objection. Protect equal participation, recognition and minority rights even when efficiency is a legitimate concern.

\textbf{Practical audit hook.} Exercise: replace three unnecessary voice votes with well-formed consent requests and identify the branch after objection.

\section{Independent review checklist}
\begin{enumerate}[leftmargin=*,itemsep=4pt]
\item Find the controlling adopted text and note its revision date.
\item Specify the decision, eligible decision maker, notice, quorum and vote basis.
\item Record evidence without fabricating approvals or preserving private member data in a public repository.
\item Classify unresolved issues as parliamentary, legal, governance policy or factual disputes.
\end{enumerate}
\end{document}
